\begin{remarks}[7.8.4]
\EGAReaderTarget{ega4.c04.item.7-8-4}{IV.2.s7.8.4-fr}
\upshape
(i) Si se consideran los anillos noetherianos $A$ que satisfacen solamente
las condiciones (ii) y (iii) de la definición (7.8.2), la inspección de las
demostraciones anteriores muestra que las afirmaciones (ii), (iii), (iv),
(v), (vi) y (vii) de (7.8.3) siguen siendo válidas si se reemplaza «excelente»
por «que satisface (7.8.2, (ii) y (iii))».

(ii) El anillo local catenario $A$ construido en (5.6.11) satisface las
condiciones (ii) y (iii) (pero no la condición (i)) de la definición (7.8.2).
En efecto, la fibra formal en el ideal maximal $\mathfrak n C_{\mathfrak n}$
de $A$ es trivialmente geométricamente regular, y la fibra en los demás
ideales primos $\mathfrak p$ de $A$ coincide con la del anillo $E$ en esos
ideales primos; ahora bien, $E$ es una $V$-álgebra de tipo finito y $V$ es
un anillo de valoración discreta, que, por tanto, es un anillo excelente
si se toma el cuerpo $k_0$ de característica $0$ (7.8.3, (ii)); por tanto,
$E$ es un anillo excelente (7.8.3, (ii)), y esto prueba que $A$ satisface
la condición (7.8.2, (ii)). Por otra parte, el complementario del punto
cerrado en $\operatorname{Spec}(A)$ es isomorfo a un abierto de
$\operatorname{Spec}(E)$, y cualquier abierto afín de este es el espectro
de un anillo
\oldpage[IV-2]{217}
excelente, que, por tanto, satisface la condición (6.12.4, $a)$) en virtud
de (7.8.3, (iv));
esto prueba que el propio $A$ satisface la condición (6.12.4, $a)$) y, por
tanto, también (6.12.4, $c)$), que no es otra que la condición (7.8.2, (iii)).
Sin embargo, $A$ no satisface la condición (7.8.2, (i)), pues hemos visto
en (5.6.11) que no es universalmente catenario.

(iii) La condición (i) de la definición (7.8.2) puede reemplazarse por la
siguiente (en apariencia más débil, teniendo en cuenta (5.6.10)):

(i$'$) \emph{Para todo ideal primo minimal} $\mathfrak p_i$ \emph{de} $A$
$(1\leq i\leq r)$, \emph{sea} $A'_i$ \emph{el cierre entero del anillo que
es un dominio de integridad} $A_i=A/\mathfrak p_i$; \emph{entonces, para
todo ideal maximal} $\mathfrak m'$ \emph{de} $A'_i$, \emph{se tiene, si
se denota por} $\mathfrak m$ \emph{la imagen inversa de} $\mathfrak m'$
\emph{en} $A$:
\begin{env}[7.8.4.1]
\label{IV.2.s7.8.4.1-fr}
\[
\dim(A_{\mathfrak m})=\dim((A'_i)_{\mathfrak m'}).
\]
\end{env}

En efecto, se ha visto (observación (i)) que, bajo las solas hipótesis
(ii) y (iii) de (7.8.2), son válidos los análogos de las propiedades
(ii), (v), (vi) y (vii) de (7.8.3). En primer lugar, se deduce de (vi) y
(ii) que los $A'_i$ satisfacen (7.8.2, (ii) y (iii)); entonces, de (v) se
deduce que los completados $((A'_i)_{\mathfrak m'})^{\widehat{\,}}$ son dominios de
integridad (e integralmente cerrados). Sea ahora $\mathfrak m$ un ideal
maximal cualquiera de $A$; para todo $i$ tal que
$\mathfrak p_i\subset\mathfrak m$, el cierre entero de
$A_{\mathfrak m}/\mathfrak p_iA_{\mathfrak m}$ es una
$(A_{\mathfrak m}/\mathfrak p_iA_{\mathfrak m})$-álgebra finita por (vi)
y, por tanto, semilocal; sus componentes locales son de la forma
$(A'_i)_{\mathfrak m'}$, donde $\mathfrak m'$ es un ideal primo
(necesariamente maximal) de $A'_i$ por encima de
$\mathfrak m/\mathfrak p_i$; de (vii), de lo anterior y de (7.8.4.1) se
deduce que los cocientes del completado
$(A_{\mathfrak m}/\mathfrak p_iA_{\mathfrak m})^{\widehat{\,}}$ por sus ideales
primos minimales tienen todos la misma dimensión, igual a
$\dim(A_{\mathfrak m})$. Por tanto (7.1.9 y 7.1.8, $b)$),
$A_{\mathfrak m}$ es formalmente catenario y, \emph{a fortiori} (7.1.11),
universalmente catenario; lo mismo ocurre, pues, con $A_{\mathfrak p}$
para todo ideal primo $\mathfrak p$ de $A$ (5.6.3, (i)), y, por tanto,
también con $A$.

En particular, si $A$ es normal o, más generalmente, si $A_{\mathfrak m}$
es unirrama para todo ideal maximal $\mathfrak m$ de $A$, la condición
(i$'$) se cumple siempre, pues $\mathfrak m$ solo puede contener un ideal
primo $\mathfrak p_i$ (\textbf 0, 16.1.5); en este caso se puede omitir (i) en la
definición (7.8.2). Más particularmente, si $A$ es un anillo noetheriano
\emph{local unirrama}, es excelente si y solo si sus fibras formales son
geométricamente regulares.

(iv) Recordemos que un anillo cociente de un anillo local de Cohen-Macaulay
(y, \emph{a fortiori}, un anillo cociente de un anillo local regular)
satisface también las propiedades (7.8.3, (ix) y (x)) en virtud de (7.2.7);
pero el interés de los cocientes de anillos regulares reside sobre todo
en sus propiedades cohomológicas (cap. III, 3\textsuperscript{a} Parte).

(v) Más adelante (\S~18) veremos que, si $k$ es un cuerpo valorado completo
no discreto, el anillo $k\{\!\{T_1,\ldots,T_n\}\!\}$ de las
\emph{series de potencias convergentes} (en una vecindad de $0$ en $k^n$)
es un \emph{anillo excelente}.
\end{remarks}
